% GENERATED FILE. Edit canonical lessons, then run npm run generate:books.
% canonical-source-hash: fnv1a64:704b2a376305eaf4
% canonical-lessons: RU-C380-predskazyvat, RU-C380-ugadyvat, RU-C380-dedlayn, RU-C380-postoyannyy, RU-C380-vremennyy

\chapter{What will happen}
\label{ch:ru-what-will-happen}
\hlchaptermodality{380}

\begin{chapteropening}
\textbf{By the end of this chapter:} \emph{I can predict what will happen, guess right, talk about a deadline, and say whether something is permanent or temporary.}

Complete the last lesson of chapter 380: I can predict what will happen, guess right, talk about a deadline, and say whether something is permanent or temporary.
\end{chapteropening}

\section[\texorpdfstring{\ru{предсказывать} (predskázyvat)}{предсказывать (predskázyvat)}]{\ru{предсказывать} (predskázyvat) — to predict, to forecast}

\label{lesson:RU-C380-predskazyvat}



\begin{quote}
Before the new one: say the Russian for to make up one’s mind, then the Russian for depending on.
\end{quote}

\subsection*{You'll want to know: \ru{предсказывать}}
\textbf{\ru{предсказывать}} (\emph{predskázyvat}) — \textquotedblleft{}to predict, to forecast\textquotedblright{}.

\begin{quote}
\ru{Трудно} \ru{предсказывать} \ru{будущее}.
\end{quote}

\begin{quote}
\emph{Trúdno predskázyvat búdushcheye.}
\end{quote}

\begin{quote}
It is hard to predict the future.
\end{quote}

Imperfective; perfective \emph{\ru{предсказать}} (\emph{\ru{предскажу}, \ru{предскажешь}}).

\begin{cousinweb}
\textbf{\ru{предсказывать}} is built from \emph{\ru{пред}-} \textquotedblleft{}before\textquotedblright{} and \emph{\ru{сказать}} \textquotedblleft{}to say\textquotedblright{}, on \emph{\ru{казать}}, Proto-Slavic \emph{*kazati} (EtymDB): to say a thing before it happens.
\end{cousinweb}

\begin{grammarlens}[title={What you've built}]
One more word for telling, explaining and describing at B1.
\end{grammarlens}

\subsection*{Your turn}
\begin{itemize}
\raggedright
  \item \emph{Say it:} \emph{predskázyvat}
  \item \emph{Say it:} \emph{predskázyvat}, once more
  \item \emph{Say it:} a sentence of your own with \emph{predskázyvat}
  \item \emph{Recall:} say the Russian for to make up one’s mind, then the Russian for depending on, then say \emph{predskázyvat} again
\end{itemize}

\begin{tcolorbox}[breakable,colback=teal!4,colframe=teal!35!black,title={Before you move on}]
What does \ru{предсказывать} mean? (\textquotedblleft{}To predict, to forecast\textquotedblright{}.) Say it once more.
\end{tcolorbox}
\section[\texorpdfstring{\ru{угадывать} (ugádyvat)}{угадывать (ugádyvat)}]{\ru{угадывать} (ugádyvat) — to guess (right)}

\label{lesson:RU-C380-ugadyvat}



\begin{quote}
Before the new one: say the Russian for depending on, then the Russian for to predict.
\end{quote}

\subsection*{You'll want to know: \ru{угадывать}}
\textbf{\ru{угадывать}} (\emph{ugádyvat}) — \textquotedblleft{}to guess (right)\textquotedblright{}.

\begin{quote}
\ru{Он} \ru{всегда} \ru{угадывает}, \ru{что} \ru{я} \ru{хочу}.
\end{quote}

\begin{quote}
\emph{On vsegdá ugádyvayet, chto ya khochú.}
\end{quote}

\begin{quote}
He always guesses what I want.
\end{quote}

Imperfective; perfective \emph{\ru{угадать}}. It means guessing right; to guess at something is \emph{\ru{догадываться}}.

\begin{cousinweb}
\textbf{\ru{угадывать}} is built from \emph{\ru{у}-} and \emph{\ru{гадать}} \textquotedblleft{}to guess, to tell fortunes\textquotedblright{}, an inherited Slavic verb: to hit on the answer.
\end{cousinweb}

\begin{grammarlens}[title={What you've built}]
One more word for telling, explaining and describing at B1.
\end{grammarlens}

\subsection*{Your turn}
\begin{itemize}
\raggedright
  \item \emph{Say it:} \emph{ugádyvat}
  \item \emph{Say it:} \emph{ugádyvat}, once more
  \item \emph{Say it:} a sentence of your own with \emph{ugádyvat}
  \item \emph{Recall:} say the Russian for depending on, then the Russian for to predict, then say \emph{ugádyvat} again
\end{itemize}

\begin{tcolorbox}[breakable,colback=teal!4,colframe=teal!35!black,title={Before you move on}]
What does \ru{угадывать} mean? (\textquotedblleft{}To guess (right)\textquotedblright{}.) Say it once more.
\end{tcolorbox}
\section[\texorpdfstring{\ru{дедлайн} (dedláyn)}{дедлайн (dedláyn)}]{\ru{дедлайн} (dedláyn) — a deadline}

\label{lesson:RU-C380-dedlayn}



\begin{quote}
Before the new one: say the Russian for to predict, then the Russian for to guess.
\end{quote}

\subsection*{You'll want to know: \ru{дедлайн}}
\textbf{\ru{дедлайн}} (\emph{dedláyn}) — \textquotedblleft{}a deadline\textquotedblright{}.

\begin{quote}
\ru{Дедлайн} — \ru{в} \ru{пятницу}.
\end{quote}

\begin{quote}
\emph{Dedláyn — v pyátnitsu.}
\end{quote}

\begin{quote}
The deadline is on Friday.
\end{quote}

Masculine. Colloquial, from office talk; in formal Russian, \emph{\ru{крайний} \ru{срок}}.

\begin{cousinweb}
\textbf{\ru{дедлайн}} is a recent borrowing from English \emph{deadline}.
\end{cousinweb}

\begin{grammarlens}[title={What you've built}]
One more word for telling, explaining and describing at B1.
\end{grammarlens}

\subsection*{Your turn}
\begin{itemize}
\raggedright
  \item \emph{Say it:} \emph{dedláyn}
  \item \emph{Say it:} \emph{dedláyn}, once more
  \item \emph{Say it:} a sentence of your own with \emph{dedláyn}
  \item \emph{Recall:} say the Russian for to predict, then the Russian for to guess, then say \emph{dedláyn} again
\end{itemize}

\begin{tcolorbox}[breakable,colback=teal!4,colframe=teal!35!black,title={Before you move on}]
What does \ru{дедлайн} mean? (\textquotedblleft{}A deadline\textquotedblright{}.) Say it once more.
\end{tcolorbox}
\section[\texorpdfstring{\ru{постоянный} (postoyánnyy)}{постоянный (postoyánnyy)}]{\ru{постоянный} (postoyánnyy) — permanent, constant}

\label{lesson:RU-C380-postoyannyy}



\begin{quote}
Before the new one: say the Russian for to guess, then the Russian for a deadline.
\end{quote}

\subsection*{You'll want to know: \ru{постоянный}}
\textbf{\ru{постоянный}} (\emph{postoyánnyy}) — \textquotedblleft{}permanent, constant\textquotedblright{}.

\begin{quote}
\ru{У} \ru{него} \ru{нет} \ru{постоянной} \ru{работы}.
\end{quote}

\begin{quote}
\emph{U negó net postoyánnoy rabóty.}
\end{quote}

\begin{quote}
He does not have a permanent job.
\end{quote}

\emph{\ru{Постоянно}} — all the time.

\begin{cousinweb}
\textbf{\ru{постоянный}} is built from \emph{\ru{по}-} and \emph{\ru{стоять}} \textquotedblleft{}to stand\textquotedblright{}, inherited from Proto-Slavic \emph{*stojati} (EtymDB): what stays standing.
\end{cousinweb}

\begin{grammarlens}[title={What you've built}]
One more word for telling, explaining and describing at B1.
\end{grammarlens}

\subsection*{Your turn}
\begin{itemize}
\raggedright
  \item \emph{Say it:} \emph{postoyánnyy}
  \item \emph{Say it:} \emph{postoyánnyy}, once more
  \item \emph{Say it:} a sentence of your own with \emph{postoyánnyy}
  \item \emph{Recall:} say the Russian for to guess, then the Russian for a deadline, then say \emph{postoyánnyy} again
\end{itemize}

\begin{tcolorbox}[breakable,colback=teal!4,colframe=teal!35!black,title={Before you move on}]
What does \ru{постоянный} mean? (\textquotedblleft{}Permanent, constant\textquotedblright{}.) Say it once more.
\end{tcolorbox}
\section[\texorpdfstring{\ru{временный} (vrémennyy)}{временный (vrémennyy)}]{\ru{временный} (vrémennyy) — temporary}

\label{lesson:RU-C380-vremennyy}



\begin{quote}
Before the new one: say the Russian for a deadline, then the Russian for permanent.
\end{quote}

\subsection*{You'll want to know: \ru{временный}}
\textbf{\ru{временный}} (\emph{vrémennyy}) — \textquotedblleft{}temporary\textquotedblright{}.

\begin{quote}
\ru{Это} \ru{временная} \ru{работа} \ru{на} \ru{лето}.
\end{quote}

\begin{quote}
\emph{Éto vrémennaya rabóta na léto.}
\end{quote}

\begin{quote}
It is a temporary job for the summer.
\end{quote}

The opposite of \emph{\ru{постоянный}}. \emph{\ru{Временно}} — for the time being.

\begin{cousinweb}
\textbf{\ru{временный}} is built on \emph{\ru{время}} \textquotedblleft{}time\textquotedblright{}, inherited from Proto-Slavic \emph{*vermę}; the \emph{-\ru{ре}-} is the Church Slavonic shape, beside native \emph{\ru{веремя}} in old texts.
\end{cousinweb}

\begin{grammarlens}[title={What you've built}]
One more word for telling, explaining and describing at B1.
\end{grammarlens}

\subsection*{Your turn}
\begin{itemize}
\raggedright
  \item \emph{Say it:} \emph{vrémennyy}
  \item \emph{Say it:} \emph{vrémennyy}, once more
  \item \emph{Say it:} a sentence of your own with \emph{vrémennyy}
  \item \emph{Recall:} say the Russian for a deadline, then the Russian for permanent, then say \emph{vrémennyy} again
\end{itemize}

\begin{tcolorbox}[breakable,colback=teal!4,colframe=teal!35!black,title={Before you move on}]
What does \ru{временный} mean? (\textquotedblleft{}Temporary\textquotedblright{}.) Say it once more.
\end{tcolorbox}
